% !TeX program = xelatex
% ============================================================================
%  第 11 课　课堂单页（学生用，单独打印一张 A4）
%  编译：XeLaTeX，两遍即可
%  提示：每一步只用上一步的余数去乘，数字永远不超过 p
% ============================================================================

\documentclass[11pt]{ctexart}
\usepackage[a4paper,left=1.7cm,right=1.7cm,top=1.4cm,bottom=1.4cm]{geometry}
\usepackage{amsmath,amssymb}
\usepackage{booktabs}
\usepackage{array}
\usepackage[dvipsnames]{xcolor}

\setCJKmainfont{SimSun}[BoldFont=SimHei,ItalicFont=KaiTi]
\IfFontExistsTF{TeX Gyre Pagella}{\setmainfont{TeX Gyre Pagella}}{\setmainfont{Latin Modern Roman}}

\definecolor{pagegreen}{RGB}{34,139,34}
\renewcommand{\arraystretch}{1.4}
\newenvironment{dongshou}[1][]{\par\vspace{3pt}}{\par\vspace{3pt}}
\newcount\xhcnt
\newcommand{\liukong}[1]{%
  \par\vspace{4pt}%
  \noindent\begin{minipage}{\linewidth}%
    \setlength{\parindent}{0pt}\setlength{\parskip}{0pt}%
    \xhcnt=#1\relax
    \loop\ifnum\xhcnt>0
      \noindent\textcolor{black!35}{\rule{\linewidth}{0.5pt}}\par\vspace{0.8cm}%
      \advance\xhcnt by -1\relax
    \repeat
    \vspace{-0.8cm}%
  \end{minipage}%
  \par\vspace{3pt}%
}
\setlength{\parindent}{0pt}
\pagestyle{empty}

\newcommand{\biaoti}[2]{%
  {\color{pagegreen}\rule{\textwidth}{1.2pt}}\\[6pt]
  {\heiti\Large 第 #1 课\quad #2}\hfill{\small\songti 课堂单页}\\[6pt]
  {\color{pagegreen}\rule{\textwidth}{0.5pt}}\\[4pt]
  {\small 姓名\underline{\hspace{3.2cm}}\qquad 班级\underline{\hspace{2.4cm}}\qquad 座号\underline{\hspace{1.6cm}}}
  \vspace{0.4em}
}

\begin{document}

\biaoti{11}{幂的尾巴转圈了：从现象到猜想}

\section*{一、填余数，找「第几步回到 1」}

\noindent 每一步只用上一步的\emph{余数}去乘。填到第一次出现 \(1\) 为止，后面的格子可以空着。

\begin{center}
{\small
\begin{tabular}{@{}l|cccccc|c@{}}
  \toprule
   & \(n=1\) & \(2\) & \(3\) & \(4\) & \(5\) & \(6\) & 第几步回到 \(1\) \\
  \midrule
  \rule{0pt}{3.2ex}\(2^n\) 除以 \(5\) & & & & & & & \\
  \rule{0pt}{3.2ex}\(3^n\) 除以 \(5\) & & & & & & & \\
  \rule{0pt}{3.2ex}\(2^n\) 除以 \(7\) & & & & & & & \\
  \rule{0pt}{3.2ex}\(3^n\) 除以 \(7\) & & & & & & & \\
  \bottomrule
\end{tabular}}
\end{center}

\vspace{0.3em}
\noindent 参考：\(2^3=8\equiv1\pmod 7\)，所以「\(2^n\) 除以 \(7\)」这一行的第一个 \(1\) 出现在 \(n=3\)。

\section*{二、算一个大的}

\begin{dongshou}
  既然余数在转圈，那么：
  \begin{itemize}
    \item \(2^n\) 除以 \(7\) 的余数，转一圈是第 \(1\) 到第 \underline{\hspace{1.4em}} 个；
    \item 所以 \(2^{100}\) 除以 \(7\) 的余数，和 \(2^{\text{第几个}}\) 一样：\\
      \(100=\) \underline{\hspace{2.4em}}，所以 \(2^{100}\equiv\) \underline{\hspace{2.4em}} \(\pmod 7\)。
  \end{itemize}
\end{dongshou}

\section*{三、把规律写下来}

\begin{dongshou}
  对着上面四个「第几步回到 \(1\)」的数字，看看它们和 \(p-1\) 有什么关系：
  \liukong{2}
  于是我猜：
  \[
    \text{当 }p\text{ 是素数、}a\text{ 不是 }p\text{ 的倍数时，}\quad a^{\,p-1}\equiv\underline{\hspace{3em}} \pmod p .
  \]
\end{dongshou}

\end{document}